\chapter{Conclusions}
In this laboratory work, I implemented and examined two Caesar ciphers for
the Romanian alphabet: a fixed-alphabet version and a version with a
keyword-based permutation. I used explicit positions from 0 to 30, rather
than character encoding values, and applied modular arithmetic to encrypt
and decrypt the text. The results confirmed that decryption with the same
key or keys recovers the prepared plaintext.

The shared normalization and validation functions handle Romanian
diacritics, cedilla variants, lowercase letters and spaces in messages.
The console retry loops reject invalid keys, unsupported characters and
keywords that are too short. The displayed permutation and mapping tables
make both algorithms easy to verify by hand. These checks helped me
understand why alphabet order and consistent preprocessing are essential
for correct encryption and decryption.

The standard version is weak because an attacker needs to try only 30
numeric keys. A keyword permutation greatly increases the theoretical
parameter space, but the second version still uses one fixed substitution
for the entire message and preserves letter frequencies and repeated
patterns. Therefore, neither cipher is suitable for modern secure
communication. The main value of these implementations is learning
substitution ciphers, modular arithmetic and reliable input handling.
